\section{Discussion}\label{sec:discussion}

The main result of this paper is the discovery of a base configuration for
$n=19$ that yields an infinite straight-line affine series attaining the optimal
number of bounded triangles for every $n=18\cdot 2^t+1$.
The base configuration was found by a systematic computational search with
two goals: to identify base configurations for infinite iteration, and to
detect arrangements that, while not supporting repeated iteration, may still
admit a single realizable iterative step.
The search was restricted to odd values of~$n$, since the iterative
constructions of~\cite{forge1998,bartholdi2007} apply only to odd~$n$;
enumeration of optimal pseudoline arrangements is also simpler in this case.

\textbf{Search for infinite series.}
The table below summarizes the results,
obtained in three stages:
enumeration of projective equivalence classes,
numerical stretchability testing, and
checking compatibility with iterative conditions.

\medskip
\noindent\begin{tabular}{|l|r|r|r|r|r|r|r|}
 \hline
 $n$                  & 15 & 17 & 19         & 21       & 23  & 25, 5-gon defect & 27 \\
 \hline
 Projective classes   & 1  & 3  & 1312       & 18       & 112 & 2324             & 56646 \\
 \hline
 Stretchable classes  & 1  & 3  & $\geq1306$ & $\geq11$ & $\geq51$ & $\geq122$   & $\geq348$   \\
 \hline
 Base configurations  & 1  & 3  & 170        & 0        & 0   & 0                & 0 \\
 \hline
\end{tabular}
\medskip

\emph{Projective classes.}
After enumerating optimal pseudoline arrangements, we identified projectively
equivalent arrangements and counted them only once.

\emph{Stretchable classes.}
Next, for each of the projective classes, we tried to realize the
corresponding pseudoline arrangements by straight lines numerically. We
minimized by gradient descent the sum of squared violations of the row-order
inequalities~\eqref{eq:general-order-ineq}. Thus, the numbers in this row are
only lower bounds: a
negative result may simply mean that the computation did not find a straight-line representative.

\emph{Base configurations.}
Finally, we tested all enumerated optimal arrangements reconstructed
from projective classes for compatibility with the conditions required
for the iterative constructions of~\cite{forge1998,bartholdi2007}.
This check is reduced to the feasibility of the linear
system~\eqref{eq:linear-order-ineq}, so the absence of a solution
can be guaranteed, unlike with the gradient-based stretching procedure.
None of the tested arrangements for $n=21$ and $n=23$
yielded a base configuration, so these zero entries
should be viewed as stronger negative evidence.
The columns $n=15$ and $n=17$ reproduce known small cases
from the series~\eqref{eq:known-series}.
The case $n=25$ is discussed separately below.

Thus, for $n=21$ and $n=23$, the
computations provide strong evidence against the existence of new
infinite iterative series of type~\cite{forge1998,bartholdi2007}.

\textbf{Infinite series for $n=25$.}
This case is more subtle because even the maximal number
$\lfloor n(n-2)/3\rfloor$ of triangles does not use all bounded segments:
when $n\equiv 1\pmod{6}$, at least two bounded segments do not belong
to any bounded triangle. In the affine plane, these \emph{unused segments} belong to
unbounded faces of optimal arrangements (see $L_{13}L_{11}L_{14}$
and $L_{12}L_{4}L_{13}$ in Figure~\ref{fig:pseudolines})
and therefore do not prevent the triangle count from attaining the upper bound.

Adding the line at infinity, however, turns the unbounded faces with those
unused segments into bounded faces whose every edge becomes unused;
we call such faces \emph{defects}.
In all optimal arrangements that we enumerated, the defect is either a
single pentagon or a pair of quadrilaterals.
For arrangements of $n$ lines in the projective plane with $n\equiv 2\pmod{6}$, the upper bound
\begin{equation}\label{eq:blanc-projective-bound}
  p_3(\cA)\le \frac{n(n-1)-5}{3}
\end{equation}
is attained for pseudoline arrangements only in the presence of a pentagonal
defect, and to our knowledge no corresponding straight-line configuration has
been reported~\cite{blanc2011}.
Accordingly, the column $n=25$ in the table concerns only pentagonal-defect
arrangements, since this is the only subclass relevant to a potentially new
projective-maximal family (arrangements with two quadrilateral defects
would only cover the same values of $n$ as the already known series
$n=6\cdot 2^t+1$~\cite{bartholdi2007}).
Within the pentagonal-defect subclass, none of the enumerated
arrangements passes the iterative compatibility test.

Note that defects introduce additional combinatorial freedom,
which explains why the counts at $n=19$ are much larger than at $n=21$
and $n=23$.

\textbf{One-step realizable examples.}
Compatibility with the iterative conditions is a stronger requirement than
ordinary stretchability, since it is designed to support repeated application
of iterative constructions. Its failure therefore rules out an infinite
iterative family of this type, but does not exclude the possibility that one
or more finite iterative steps still produce straight-line realizable
arrangements. Explicit realizable arrangements of
$41$ and $45$ lines, obtained from arrangements of $21$ and $23$ lines after a
single iterative step, are recorded in Appendix~\ref{sec:appendix-one-step}.
Besides these examples, we found a realizable arrangement
of $49$ lines obtained from a $25$-line arrangement with a pentagonal defect
after a single iterative step. After adding the line at infinity, this yields
a projective arrangement with $50$ lines and $815$ triangles, attaining the
projective upper bound \eqref{eq:blanc-projective-bound}.

Finally, we observe that as $n$ increases, the proportion of stretchable
arrangements among the enumerated optimal pseudoline arrangements appears to
decrease. We conjecture that the number of infinite series of
type~\cite{forge1998,bartholdi2007} arising from stretchable
optimal arrangements is finite.
